\documentclass[11pt,a4paper]{article}
\usepackage[T1]{fontenc}
\usepackage[utf8]{inputenc}
\usepackage{lmodern,microtype}
\usepackage[margin=25mm]{geometry}
\usepackage{amsmath,amssymb,amsthm,mathtools,mathrsfs}
\usepackage{booktabs,longtable,array,enumitem}
\usepackage[dvipsnames]{xcolor}
\usepackage{graphicx}
\usepackage{url}
\usepackage[unicode,colorlinks=true,linkcolor=MidnightBlue,citecolor=MidnightBlue,urlcolor=MidnightBlue]{hyperref}
\usepackage{fancyhdr}
\pagestyle{fancy}
\fancyhf{}
\fancyhead[L]{\small Explicit Diophantine forms for $\operatorname{Con}(\mathrm{ZFC})$}
\fancyhead[R]{\small ProveIt research supplement}
\fancyfoot[C]{\thepage}
\setlength{\headheight}{14pt}
\setlength{\emergencystretch}{2em}
\setlength{\parskip}{3pt}
\setlist{nosep}
\allowdisplaybreaks[2]
\numberwithin{equation}{section}
\newcommand{\N}{\mathbb N}
\newcommand{\Z}{\mathbb Z}
\newcommand{\Q}{\mathbb Q}
\newcommand{\Con}{\operatorname{Con}}
\newcommand{\Sq}{\operatorname{Sq}}
\newcommand{\ZFC}{\mathrm{ZFC}}
\newtheorem{theorem}{Theorem}[section]
\newtheorem{lemma}[theorem]{Lemma}
\newtheorem{proposition}[theorem]{Proposition}
\newtheorem{corollary}[theorem]{Corollary}
\theoremstyle{definition}
\newtheorem{definition}[theorem]{Definition}
\newtheorem{remark}[theorem]{Remark}
\title{Smaller explicit Diophantine forms for\newline $\Con(\ZFC)$\\[6pt]
\large Ten integer unknowns, compact quadratic systems,\newline and a smaller fixed coefficient instance}
\author{Research synthesis prepared for the ProveIt repository}
\date{October 2026\\\small Source snapshot \texttt{8c6eec9d19ef80538c8e61ff6e0180040583446a}}

\begin{document}
\maketitle
\begin{abstract}
We give explicit polynomial equations whose absence of solutions is
provably equivalent in ZFC to its syntactic consistency.  The fixed
instance is derived from ProveIt's calibrated blank-input K215 machine;
no unspecified recursively enumerable index or existential coefficient
is used.  A specialization of Sun's ten-integer construction, together
with a coprime square merger and a quintic divisibility merger, gives an
equation in ten unrestricted integer variables.  A separately developed
Jones-based equation uses sixteen unrestricted integer variables and
1,617 binary arithmetic operations.  For readers prioritizing degree or
evaluation cost, we give a system of 22 quadratic equations in 31 positive
variables evaluated in 60 operations, its 125-operation quartic sum of
squares, and a 12-positive-variable system evaluated in 102 operations.
All these compatible forms can use a new complete native coefficient
tuple whose largest port has 1,070,004 bits, 45.36\% below the compared
repository instance.  The archive contains ordinary mathematical proofs,
literal coefficient data, division-free circuits and reproducible finite
audits.  It does not claim global optimality or a new end-to-end
proof-assistant certificate.
\end{abstract}
\tableofcontents
\clearpage

\section{The target and the meaning of ``simpler''}
\label{sec:overview}

Let $\Con(\ZFC)$ mean that there is no finite derivation of a contradiction
from the recursively presented first-order ZFC axioms, in the proof
calculus used by the repository's K215 calibration.  Equivalent standard
effective proof calculi give provably equivalent consistency statements.
The required conclusion has the following form for one fully fixed
$F\in\Z[X_1,\ldots,X_n]$:
\begin{equation}\label{eq:target}
 \ZFC\ \vdash\quad
 \bigl[\neg\exists\boldsymbol{x}\in\Z^n\ F(\boldsymbol{x})=0\bigr]
 \ \longleftrightarrow\ \Con(\ZFC).
\end{equation}
This is a proof of a biconditional.  It does not prove that either side
holds.  Under the usual consistency hypothesis, G\"odel's second
incompleteness theorem precludes ZFC from proving the left side once
the biconditional has been established.

There are several inequivalent notions of simplicity.  Fewer variables
can cost enormous degree; a low-degree system can require more witnesses;
small arithmetic circuits can contain very large fixed integers.  A
positive-variable equation also differs from an equation in unrestricted
integers.  We therefore retain several useful points instead of calling
one construction an absolute minimum.

\subsection{Concrete improvements over the pinned collection}

All comparisons are against the immutable snapshot in \cite{repo}, not
against a moving branch.  ``Operation'' means a binary addition,
subtraction or multiplication in the emitted directed acyclic graph.
Sharing is allowed; all powers of variables are expanded into such
gates.  Fixed integers are leaves.  A system is evaluated as shared left
and right sides to be compared; the comparisons themselves are not
arithmetic gates.  A sum of squared residuals includes the subtraction,
squaring and summation gates.  These conventions match the compared
repository receipts.

\begin{center}
\small
\begin{tabular}{@{}p{43mm}rrrr@{}}
\toprule
Form & Domain & Variables & Degree & Operations\\
\midrule
New Sun/Jones synthesis & $\Z$ & 10 & $d_{10}$ & 1,442\\
Jones wrapper, package 29 & $\Z$ & 17 & $d_{17}$ & 20,607\\
New Jones wrapper & $\Z$ & 16 & $d_{16}$ & 1,617\\
Quadratic system, package 30 & $\Z_{>0}$ & 31 & 2 & 65\\
New quadratic system & $\Z_{>0}$ & 31 & 2 & 60\\
Quartic, package 30 & $\Z_{>0}$ & 31 & 4 & 130\\
New quartic & $\Z_{>0}$ & 31 & 4 & 125\\
Product, package 30 & $\Z_{>0}$ & 12 & 1,081 & 111\\
New product evaluator & $\Z_{>0}$ & 12 & 1,081 & 108\\
System, package 30 & $\Z_{>0}$ & 12 & 1,072 & 107\\
New shifted system & $\Z_{>0}$ & 12 & 1,072 & 102\\
\bottomrule
\end{tabular}
\end{center}
Here
\[
 d_{10}=126496\,5^{31}+27064,\qquad
\]
\[
 d_{17}=424296\,5^{33}+53824,\qquad
 d_{16}=813056\,5^{32}+10016.
\]
The new ten-integer equation is specified in full in
Section~\ref{sec:signed10}.  Its exact degree is
\[
589042901992797851562527064,
\]
and its evaluator uses 851 multiplications
and 591 additions/subtractions.  It combines the fixed quartic input
with Sun's ten-integer reduction \cite{sun2021,sun2026} and the local
simplifications proved there.  The contribution is an explicit
$\Con(\ZFC)$ instance and smaller local formulas, not a new general
ten-unknowns theorem.

The collection already contains nine-variable equations over
$\N$ or $\Z_{>0}$.  Ten unrestricted integer variables should not be
compared with these as if the domains were identical.  The existing
67-operation single equation and the 59-operation higher-degree system
also remain better at their respective operation-count objectives.
Our 60-operation system improves the quadratic/quartic point.

\subsection{Fixed coefficients and the proof chain}

All the new forms are tied to one complete native coefficient tuple.
For its four ports the largest bit length falls from 1,958,376 to
1,070,004, while their combined raw storage falls from 977,220 to
533,099 bytes.  These are measurements of the \emph{native ports}; they
are not a claim about the largest coefficient after expanding an entire
high-degree polynomial.  In particular, the sixteen-variable Jones
wrapper has a much larger fixed base coefficient specified by a bounded
power-of-two recipe.

The logical composition is
\begin{align*}
 \text{a zero of the displayed equation/system}
 &\ \longleftrightarrow\ \text{a positive zero of the native parent}\\
 &\ \longleftrightarrow\ \text{a finite blank-input K215 halt}\\
 &\ \longleftrightarrow\ \text{a finite ZFC refutation}.
\end{align*}
Negating this equivalence yields \eqref{eq:target}.  The parent-to-machine
direction uses the native compiler and arithmetic-skeleton proofs in
the pinned repository.  The last equivalence is calibrated by
\texttt{blank\_halts\_iff\_zfc\_refutation} and
\texttt{nonhalting\_iff\_zfc\_consistency} in
\texttt{KEndpoint215.lean}.  The replacement tuple is justified by the
ordinary adapter proof in Section~\ref{sec:new-coefficients}.

No coefficient depends on deciding whether K215 halts.  The 430 machine
transitions, 943 tiles, masks, weighted quartic coefficients and all fixed
powers are bounded, explicitly reconstructible data.  Auxiliary names
inside the displayed formulas denote polynomial expressions, not
additional existential variables.  Fixed sums with large numerical
limits are ordinary polynomials; their efficient evaluators are supplied.

\subsection{Verification boundary}

The archive separates three kinds of evidence.  First are the inherited
semantic results, pinned source proofs and the existing Lean endpoint.
Second are the new ordinary mathematical proofs: the coefficient adapter,
square and divisibility mergers, the cofinal Jones restriction, and the
algebraic refactorings.  Third are executable finite checks of coefficient
bytes, symbolic identities, degree witnesses and circuit transcription.
The finite checks are useful for catching errors in a concrete equation;
they do not prove universal compiler correctness by testing examples.

The new work is not an end-to-end Lean or Rocq development.  In particular,
the ten-variable route imports a theorem from the explicitly identified
2026 preprint \cite{sun2026}; the sixteen-variable route independently
uses the Jones reduction and does not depend on that preprint.  The
ordinary proofs use finite arithmetic, induction, elementary algebra and
the cited Diophantine theorems, all available within ZFC.  No truth or
consistency assumption about ZFC is needed to prove the biconditional.

\section{A smaller complete coefficient tuple}
\label{sec:new-coefficients}

The mixed-radix clause encoding of package 25 and the arbitrary-odd-factor
rank controller of package 24 combine without changing the 943-tile
language.  This removes two independent power-of-five roundings:
\[
 \rho=2^{121},\qquad L=8801,\qquad
 d=121\cdot8801=1064921,\qquad B=2^d.
\]
The largest fixed port has 1,070,004 bits, compared with 1,958,376 bits
in packages 25 and 27.  The four raw payloads occupy 533,099 bytes instead
of 977,220 bytes.  These reductions are approximately $45.36\%$ and
$45.45\%$, respectively.  Variables, degree and arithmetic gate count of
a compatible skeleton do not change.

\subsection{The literal bounded recipe}

Retain the ordered 943 tiles in the bundled literal table, in slot order
west, north, east, south.  Start is tile 0 and the impossible tile is
tile 1.  The new builder reconstructs this ordered table from all 430
literal K215 transitions.  Put
\begin{align*}
 \mathcal E={}&\{0,\ldots,942\}\cup\{983,\ldots,992\}
               \cup\{1003,\ldots,1032\}\\
 &{}\cup\{1063,1064,1095,1096,2195,3296\}.
\end{align*}
The four copy offsets are $1023,1013,1003,983$; the four anchors are
$1095,1096,2195,3296$.  Set $W_j=2^{10+2j}$ for $0\le j<44$.
If tile $i$ has edge labels $x_{it}$, define
\[
 c_i=1+\sum_{t=0}^3\sum_{j=0}^9
       \operatorname{bit}_j(x_{it}+512)W_{10t+j}
       +\sum_{j=40}^{43}W_j.
\]
At copy exponent $\operatorname{offset}_t+j$, set $c_e=W_{10t+j}$.
The four anchors receive $W_{40},\ldots,W_{43}$ in the listed order;
$c_{1063}=c_{1064}=0$.  Put
\[
 \mu=1022+\sum_{j=0}^{43}W_j,\qquad
 Z=(2^{119}-1)\sum_{e=8792}^{8798}\rho^e
       +(2^{27}-1)\rho^{8799}.
\]
The exact constants are
\begin{align*}
 D_C={}&\rho^{30}+\rho^{62}+\rho^{1100}+\rho^{2200}
       +\sum_{e\in\mathcal E}c_e
            (\rho^{4393-e}+\rho^{5495-e}),\\
 D_R={}&\rho^{42},\\
 M_C={}&B-1-\sum_{\substack{e\in\mathcal E\\e\ne1}}\rho^e
                -2\rho^{1063},\\
 M_{F,0}={}&4+\mu(\rho^{4393}+\rho^{5495})
       +\sum_{j=0}^{9}(\rho^{1013+j}+\rho^{1065+j})
       +\rho^{2195}+\rho^{3296}+Z.
\end{align*}
The four ports, always replaced together, are
\begin{equation}\label{eq:new-four-ports}
 \boxed{\ \beta=B-1,\quad K=D_C+BD_R,\quad
 M_C,\quad M_F=M_{F,0}+B-1.\ }
\end{equation}
The last port is already shifted.  Every number in this definition is
computed by a bounded finite recipe, with no search for a halting run.

\subsection{What the finite geometry certifies}

The clause coefficients and copy/anchor geometry are unchanged from
package 25.  If $n$ selectors are active, then $0\le n\le943$ and the low
ten bits of the clause sum equal $n$.  The mask 1022 forces $n=0$ or 1,
without any assumption about carry freedom in higher fields.  Once this
is known, every radix-four field is a sum of two Boolean values.  Its
masked low bit forces the copy or anchor to equal the selected value.

The new cell circumference satisfies the exact identities
\begin{align*}
 \{(4393-e)\bmod8801:e\in\mathcal E\}
 &\cap\{(5495-e)\bmod8801:e\in\mathcal E\}=\varnothing,\\
 \{(2195-e)\bmod8801:e\in\mathcal E\}
 &\cap\{(3296-e)\bmod8801:e\in\mathcal E\}=\{0\}.
\end{align*}
Thus there is a clean clause band under every rotation, and the two
occupied anchors force whole-cell alignment.  The ordinary overlap
identities are the same as in package 25: targets $1013+j$ compare
sources $983+j$ and $1013+j$, while targets $1065+j$ compare $1003+j$
and $1023+j$; the anchor targets compare $1095$ with $2195$ and $1096$
with $3296$.  The builder checks all 69,230 scalar basis identities.
Linearity extends them to every native assignment.

The uniform pre-decoding bound is
\[
 \mathcal M=(989+2)(2\sum_e c_e+6),\qquad
 T_0=\max(4\mathcal M+8,2\mu+4,16).
\]
Exact arithmetic gives $2^{119}<T_0<2^{120}<\rho$.  Hence 121 is the least
odd radix width for this unchanged bound.  This is not a global minimum.
The maximum unshifted product exponent is $5495+3296=8791$, so all eight
padding cells are identically zero in canonical assignments, including
after every allowed dummy change.  Their additional tests preserve
completeness and may be ignored in soundness.  Also
$2(2^{119}-1)+4<\rho$, so the mask fits the inherited slack estimate.

The padding contributes $7\cdot119+27=860$ bits; the core high mask has
$2\cdot53+20+2+1=129$ bits.  Materialization confirms
\begin{gather*}
 \operatorname{popcount}(M_{F,0})=989,\qquad
 \operatorname{popcount}(M_C)=d-989,\\
 0<M_C,M_{F,0}<B-1,\qquad
 M_C\equiv2\pmod4,\qquad M_{F,0}\equiv4\pmod8.
\end{gather*}

\subsection{An explicit controller for the actual packed rank}

The new $d$ is not a power of five, so the old five-adic tuning statement
cannot simply be reused.  Factor $d=11^2\cdot13\cdot677$ and set
\[
 \lambda=2341,\qquad k=2340=\lambda-1.
\]
Trial division proves that $\lambda$ is prime.  Exact modular arithmetic
gives the following orders:
\begin{center}
\begin{tabular}{c|rrrr}
Modulus & 121 & 13 & 677 & 2341\\ \hline
Order of $B$ & 10 & 12 & 52 & 60
\end{tabular}
\end{center}
Thus $\operatorname{ord}_d(B)=780\mid k$, and
\[
 B^{2340}\bmod2341^2=2556373=1+1092\cdot2341,
 \qquad v_\lambda(B^k-1)=1.
\]
For $H_0=1+BD_C+B^2(D_R+1)$ the necessary unit residues are
\begin{center}
\begin{tabular}{r|rrr}
Prime $p$ & $B\bmod p$ & $(B^2-1)\bmod p$ & $H_0\bmod p$\\ \hline
11   & 2    & 3    & 9\\
13   & 6    & 9    & 7\\
677  & 611  & 293  & 654\\
2341 & 2100 & 1896 & 1225
\end{tabular}
\end{center}

Choose sufficiently large independent powers of $\lambda$ for tableau
width $w$ and height $h$, and put $N=wh=\lambda^n$, $q=B^N$.
The inherited language allows arbitrary sufficiently large independent
rectangular padding.  Impose $N>d\lambda^2$ and $h>\lambda$.
The order table gives $q\equiv B^w\equiv B\pmod{d\lambda}$.
Activating the lower dummy in cell $j$ changes the actual packed Pell
index by the inherited exact formula
\[
 \delta R=-\Gamma B^j,\qquad
 \Gamma=\rho^{1063}(q^2-1)
             [1+q(D_C+BD_R+B^w)].
\]
The unit table proves $\gcd(\Gamma,dN)=1$.

\paragraph{Boolean coverage.}
Set $T=N/\lambda$ and $C_0=d\lambda$.  Every residue modulo $dN=C_0T$
is a Boolean subset sum of
\[
 \{B^{2+kj}:0\le j<T\}\cup\{B^3\}.
\]
Indeed, binomial lifting and $B^k\equiv1\pmod d$ show that its first
$T$ powers are exactly the residues $1+C_0z$, $z$ modulo $T$.
Normalize the desired residue by $B^{-2}$ and choose
$\varepsilon\in\{0,1\}$ so that subtracting $\varepsilon B$ leaves a
nonzero residue modulo $\lambda$.  Let $a$ be the least residue modulo
$C_0$.  Then $1\le a<C_0<T$ and $\gcd(a,T)=1$.  Write the target as
$a+C_0u$, solve
\[
 az_0+\frac{a(a-1)}2\equiv u\pmod T,
\]
and select the $a$ distinct orbit entries whose additive coordinates are
$z_0,\ldots,z_0+a-1$.  Include $B^3$ exactly when $\varepsilon=1$.
This constructs the required subset with no repeated position.
The largest cell index is
\[
 2+(\lambda-1)(N/\lambda-1)=N-N/\lambda-\lambda+3.
\]
Since $w<N/\lambda$, its translates by one and by $w$ do not wrap.
The extra position 3 is distinct because $k>1$.

Start with zero control bits and actual rank $R_0$.  Applying coverage to
$(R_0-dw)\Gamma^{-1}$ tunes $R\equiv dw\pmod{dN}$, exactly the temporal
rotation required by $X=2^R$.  The previously proved carry, slack and
padding bounds hold under every selected change.  This proves unbounded
completeness; the finite modular certificates alone would not do so.

\subsection{Full interface and verification boundary}

All of $b,L,d,N$ are odd.  With $t=dN$, the canonical Start and low-mask
congruences still give $R\equiv3\pmod4$, while the exact mask population
identity gives $\operatorname{popcount}(R)=3t+2$.  Therefore
$X/q=2^{R-t}$ is a positive square.  The inherited even main-Pell
parameter and strong auxiliary index preserve the odd positive multiplier;
the half-binomial valuation gives $v_2(Y)=3t$, so $s=Y/q^3$ remains odd.
These are precisely the extra completeness promises needed for the
compatible signed reductions.

Consequently the complete tuple \eqref{eq:new-four-ports} satisfies the
same blank-K215 native interface as the source tuple, including its
positive slack and signed parity requirements.  Composing this ordinary
adapter proof with a compatible arithmetic-skeleton proof and the K215
endpoint preserves the ZFC-provable equivalence between solvability and
existence of a finite ZFC refutation.

The bundled standard-library builder reconstructs the exact payloads,
tests all finite geometry identities and modular certificates, and records
hashes and bit lengths.  It does not decide consistency or supply an
end-to-end Lean/Rocq certificate.  The complete new adapter proof is also
provided separately, with the relevant pinned source proofs and notices.

\section{The common positive quartic and its literal index}
\label{sec:parent}

We use the unshifted positive quartic of package 30 as the coding source.
This choice keeps the inherited coefficient index unchanged as a
symbolic recipe.  The cheaper 60-operation system later in this report
is an equivalent presentation; it is not silently substituted into the
index compiler.

Let $\beta,K_0,M_C,M_F$ denote the complete tuple of
\eqref{eq:new-four-ports}; $K_0$ is that equation's port $K$ and is distinct
from subsequent polynomial expressions also traditionally called $K$.
The 31 positive coordinates, in their fixed order, are
\[
 (J,F,\alpha,z,f,h,i,o,j,s,w,\tau,\eta,\zeta,y,\gamma,
 N,E,P,c,a,X,D,C_2,L,F_2,V,g,I,B_{\mathrm{aux}},\Delta).
\]
Put
\[
 q=\beta J+1,\quad k=\eta+\zeta,\quad
 H=4a+3,\quad C=q-F-\alpha.
\]
The parent is the sum of the squares of the residuals of these 22
equations, with ``left minus right'' as the residual convention:
\begin{alignat}{2}
 L&=q^2-1, &\qquad N&=q(L+1),\notag\\
 X&=wq, & a&=E+sN,\notag\\
 E&=X(a-E), & c&=k(a-E)+\eta,\notag\\
 P&=E(c-\eta)+k, & \tau^2&=P(P-k)+1,\notag\\
 D&=X+ac+\gamma H, & \Delta&=a^2+H,\notag\\
 C_2&=c^2, & D^2&=\Delta C_2+1,\notag\\
 k&=jc-V+1+hE, & F_2&=f^2,\notag\\
 V&=of-c, & g&=y^2-V^2,\notag\\
 I&=iC_2, & B_{\mathrm{aux}}&=I^2,\notag\\
 B_{\mathrm{aux}}&=\Delta(F_2-1), & y^2-B_{\mathrm{aux}}g&=1,
 \label{eq:parent-quadratic}
\end{alignat}
together with the two transport equations
\begin{align}
 jc-V&=(q-F)(N-q)-CL+(M_C+qM_F)J,\label{eq:parent-mask}\\
 (K_0+w)C+q-F-z(q-1)&=1.\label{eq:parent-transport}
\end{align}
The order in the literal JSON, rather than the grouping on this page,
is authoritative for the coefficient indices.  There are no products of
more than two coordinates in a residual after the displayed abbreviations
are substituted: every residual has total degree at most two.

Write this sum of squares as $Q_{31}$ and let
\[
 P_\star(x_0,x_1,\ldots,x_{31})
   =Q_{31}(x_1+1,\ldots,x_{31}+1).
\]
The dummy input $x_0$ does not occur.  The coefficient of $\tau^4$ in
$Q_{31}$ is 1, so the degree is exactly four.  At the all-zero
nonnegative input the residual $E-X(a-E)$ is 1, hence
$P_\star(a,0,\ldots,0)\geq1$ for every integer $a$.  Thus Sun's source
construction applies directly; there is no need to append a variable or
square this quartic again.

\begin{proposition}\label{prop:parent}
With the complete tuple \eqref{eq:new-four-ports}, ZFC proves
\[
 (\exists\boldsymbol{n}\in\N^{31}\ P_\star(0,\boldsymbol{n})=0)
 \ \longleftrightarrow\ 
 (\text{there is a finite ZFC refutation}).
\]
\end{proposition}
\begin{proof}
Translation by 1 is a bijection from $\N$ to $\Z_{>0}$.  A sum of integer
squares is zero exactly when every residual is zero.  The package-30
parent proof identifies these positive solutions with the compatible
native arithmetic skeleton.  Section~\ref{sec:new-coefficients} supplies
that skeleton with a complete tuple preserving its soundness,
completeness, slack and parity interface.  The native compiler and the
K215 endpoint then give the assertion.  The strong auxiliary Pell index
must be converted to the ordinary index as in the inherited parent
proof; the tuple replacement does not change that conversion.
\end{proof}

\subsection{Weighted coefficients with no unspecified index}

Put $\nu=31$ and $L_\star=5^{32}=23283064365386962890625$.  Write
\[
 P_\star=\sum_{|\boldsymbol{a}|\le4}p_{\boldsymbol a}
                x_0^{a_0}\cdots x_{31}^{a_{31}},\qquad
 c_{\boldsymbol a}=(4-|\boldsymbol a|)!
              \prod_{j=0}^{31}a_j!\ p_{\boldsymbol a}.
\]
Define $e(\boldsymbol a)=\sum_{j=0}^{31}a_j5^j$.
All $58,905=\binom{36}{4}$ positions of total weight at most four are
distinct, and $e(\boldsymbol a)<L_\star$.  There are 402 nonzero weighted
coefficient records when the four ports are treated as fixed symbolic
parameters.  The file \texttt{signed16/data/constants.json} gives all of
them as explicit finite recipes.  The verifier expands the normalized
parent symbolically and checks every record, including factorial weights.

For the new tuple, the explicitly materialized integer
\[
 z_\star=2^{4269855}
\]
satisfies
\begin{equation}\label{eq:coefficient-bound}
 |c_{\boldsymbol a}|+1<z_\star
 \quad\text{for every }|\boldsymbol a|\le4.
\end{equation}
The check uses exact integer arithmetic on the actual port payloads,
not a floating-point estimate.  This also bounds every
$|p_{\boldsymbol a}|$ by $z_\star$, since its factorial multiplier is a
positive integer.

For the Jones route define the fixed polynomial and fixed integers
\begin{align*}
 e_\star(T)&=\sum_{|\boldsymbol a|\le4}
                  (z_\star+c_{\boldsymbol a})T^{L_\star-e(\boldsymbol a)},\\
 y_\star&=e_\star(2z_\star),\qquad
 \kappa=(2z_\star)^{L_\star+1}
        =2^{99415332078933715820316769856}.
\end{align*}
Every digit lies strictly between 0 and $2z_\star$ and its position is
distinct, so $0<y_\star<\kappa$.  These are bounded definitions of actual
integers.  The huge $y_\star$ need not be a leaf of a circuit, and the
power-of-two recipe for $\kappa$ is explicit even though expanding its
binary string is not practical.

For the ten-variable route we instead use the sparse polynomial
\[
 D_\star(T)=\sum_{|\boldsymbol a|\le4}
                   c_{\boldsymbol a}T^{L_\star-e(\boldsymbol a)}.
\]
It has exactly the same weighted coefficient data, without the all-position
$z_\star$ padding.  This is one reason to preserve the literal parent and
its coordinate order across the two independent reductions.

\section{An explicit equation in ten unrestricted integers}
\label{sec:signed10}

The smallest unknown count in this package is achieved by the following
fully fixed equation:
\[
 F_{10}(f,g,h,k,l,w,x,y,z,m)=0,
 \qquad(f,g,h,k,l,w,x,y,z,m)\in\mathbb Z^{10}.
\]
Its division-free circuit uses \textbf{1,442 binary arithmetic operations}
and has exact total degree
\[
 126496\cdot5^{31}+27064
 =589042901992797851562527064.
\]
The general ten-integer reduction is due to Sun's 2026 preprint; this
report does not claim a new universal unknown-count theorem.  Our
contribution here is its completely specified Con(ZFC) specialization,
together with three local reductions of its degree and arithmetic circuit.
The two primary sources are Sun's 2021 paper, especially the explicit
proof of Theorem~3.1, and his 2026 preprint, Sections~3, 6 and~7.\footnote{
Z.-W.~Sun, \emph{Further results on Hilbert's Tenth Problem},
Sci. China Math. 64 (2021), 281--306,
\url{https://arxiv.org/html/1704.03504v7};
Z.-W.~Sun, \emph{Ten unknowns for Hilbert's tenth problem over the integers},
2026 preprint, \url{https://arxiv.org/html/2609.01594v2}.}
The literal Sun final formula is retained as a reference variant.

\subsection{Fixing every source coefficient}

Let $r_1,\ldots,r_{22}$ denote the quadratic residuals of the positive
31-variable system, with the complete new four-port tuple.  Define
\[
 P_0(n_1,\ldots,n_{31})
 =\sum_{j=1}^{22}r_j(n_1+1,\ldots,n_{31}+1)^2.
\]
The preceding reduction proves that a nonnegative zero of $P_0$ exists
exactly when ZFC is inconsistent.  Its degree is four, and $P_0(0)>0$:
the fifth residual $E-X(a-E)$ equals one at the all-one parent assignment.
Apply Sun's source construction to
\[
 P(a,n_1,\ldots,n_{31})=P_0(n_1,\ldots,n_{31}),\qquad a=0.
\]
The input $a$ is dummy.  This direct specialization avoids any unspecified
index for a recursively enumerable set.

For $\alpha\in\mathbb N^{32}$ with $|\alpha|\le4$, write $p_\alpha$ for
the corresponding ordinary coefficient of $P$ and set
\[
 c_\alpha=(4-|\alpha|)!\prod_{j=0}^{31}\alpha_j!\,p_\alpha,
 \qquad e(\alpha)=\sum_{j=0}^{31}\alpha_j5^j.
\]
Exactly 402 nonzero weighted coefficient recipes are retained in
\texttt{signed10/data/constants.json}; unlisted coefficients are zero.
The verifier reconstructs all of them from the shifted residual squares.
All nonzero records have $\alpha_0=0$.  The checked coefficient bound is
\[
 |c_\alpha|+1<2^{4269855}.
\]
Fix
\[
 \beta_S=2^{4269882}=2^{2\cdot2134941}.
\]
Because $48\cdot33^4=56924208<2^{27}$ and $|p_\alpha|\le|c_\alpha|$,
\[
 \beta_S>33^4\,4!\,2\max_\alpha|p_\alpha|.
\]
This is precisely the source bound in Sun 2021 for $\nu=31$, $\delta=4$
and $p=2$.  The additional constant is supplied both by its exact power
and by a compressed binary payload.  The binding of every fixed leaf is
\texttt{signed10/data/INSTANCE.json}.

\subsection{The complete source polynomials}

Put $t=5^{31}$ and $L_c=5^{32}=5t$.  Define
\begin{align*}
 b&=1+3f,&\mathcal B&=\beta_S b^4,\\
 \ell&=\sum_{j=1}^{31}\mathcal B^{5^j},&
 D_c&=\sum_{|\alpha|\le4}c_\alpha\mathcal B^{L_c-e(\alpha)},\\
 N_0&=\mathcal B^{t+1},&N_1&=4\mathcal B^{9t+1},\\
 N&=N_0N_1,&M_c&=N_0-1-(b-1)\ell,\\
 T_c&=M_c+(\mathcal B-2)\mathcal B^{L_c}N_0,&
 J_c&=2(1+g)^4D_c+\sum_{i=0}^{9t}\mathcal B^{i+1},\\
 S_c&=g+J_cN_0,&R_c&=(S_c+T_c+1)N+T_c+1,\\
 X&=(N-1)R_c,&Y&=N^2,\\
 \mathscr C&=b\mathcal B^t.&&
\end{align*}
All powers and summation limits are fixed integers.  The apparent division
by $p-1$ in Sun's general formula disappears at $p=2$.  The mask $M_c$
is his digit mask in telescoped form.  The bound has the prescribed form
$\mathscr C=2^{4269882t}b^{4t+1}$.

Sun's source theorem supplies, from a nonnegative zero of $P_0$, witnesses
with $f>0$, $b$ a square and a power of two, $b\le g<\mathscr C$, and
$Y\mid\binom{2X}{X}$.  Conversely, $f\ne0$, $0\le g<2\mathscr C$, the
same square and power conditions on $b$, and that divisibility imply a
nonnegative zero of $P_0$.  The construction also gives the strict size
bounds $X>3b$ and $Y>\max\{b,256\}$ required by Sun 2026.  Thus all source
hypotheses are checked for this particular, fully fixed polynomial.

\subsection{The bridge and the shared Lucas witnesses}

Use the following polynomial abbreviations, in the displayed order:
\begin{align*}
 L&=lY,& U&=2LX,& V&=4gwY,& W&=bw,\\
 A&=U(V+1),& B&=2X+1,& C&=B+(A-2)h,\\
 P_n&=U^2V,& K&=X+1+k(P_n-2),\\
 S&=2A-5,&T&=3WC-2(W^2-1),\\
 R_2&=(P_n^2-4)K^2+4.&&
\end{align*}
Then set
\begin{align*}
 \Delta&=8\mathscr C^3gK^2
 -g^2\bigl(32(C-KL)^2\mathscr C^3+g^2K^2\bigr),\\
 O&=f^2l^2\Delta,\qquad \xi=zx,\\
 D&=(A^2-4)C^2+4,\qquad E=C^2D\xi,\\
 F&=4(A^2-4)E^2+1,\\
 G^*&=2(A+2)(A-2)^2E^2-1,\qquad \Omega=16E,\\
 H^*&=C+(B-C)F+\Omega Fy,\\
 I^*&=((G^*)^2-1)(H^*)^2+1.
\end{align*}
These are Sun 2026's formulas (6.1)--(6.7), followed by his modified Lucas
certificate (3.1)--(3.7).  The symbol $P_n$ only factors a repeated product.
The integer $S$ is always nonzero and odd.

His positive bridge supplies $b$ square, $R_2$ square, $S\mid T$,
$C=u_B(A)$, and $O>0$.  In fact $O\ge4$: on this branch $g\ge b\ge4$,
and $\Delta>0$ is a multiple of $g$; also $f,l\ne0$.
His reverse bridge states that
\[
 b\text{ square},\quad O>0,\quad\xi\ne0,\quad
 DFI^*\text{ square},\quad R_2\text{ square},\quad S\mid T
\]
imply a nonnegative zero of $P_0$.  These precise bridge statements are
the external Lucas-theoretic inputs used below.

\subsection{Three local simplifications}

\paragraph{A quintic merger.}
Use
\[
 \pi(v)=(v^3-v-1)(v^2+23),\qquad
 \Sigma=\xi S^5,
\]
and
\[
 r=T+Sy,\qquad
 \Theta=(r^3-rS^2-S^3)(r^2+23S^2).
\]
The polynomial $\pi$ has a root modulo every positive integer.  For an
odd prime other than 23, either its cubic factor has a simple root, or
that cubic is irreducible and its discriminant $-23$ is a square: the
Frobenius permutation is an even 3-cycle.  In the latter case the quadratic
factor has a simple root.  At 23 the cubic has the simple root 3.  At 2,
$-23\equiv1\pmod8$ has a 2-adic square root.  Hensel lifting and the Chinese
remainder theorem finish the argument.  Neither factor has a rational root.

If $S\mid T$ and $\xi\ne0$, prescribe a residue class of $y$ modulo $\xi$
so that $\Sigma\mid\Theta$; Sun's full residue interpolation for the Lucas
witness allows exactly this prescription.  Conversely, if a prime has
$\nu_p(S)>\nu_p(T)$, the two factors of $\Theta$ have valuations
$3\nu_p(T)$ and $2\nu_p(T)$, contradicting $S^5\mid\Theta$.  Hence
$\Sigma\mid\Theta$ forces $S\mid T$.  Moreover $\Theta\ne0$ whenever
$S\ne0$, by the absence of rational roots.  This replaces Sun's degree-six
merger with a degree-five one.

\paragraph{One square in place of two.}
Define
\[
 r_0=b\bigl(1+(3O-4)b^2z^2\bigr),\qquad
 r_1=R_2,\qquad r_2=DFI^*.
\]
For $z\ne0$, a square $r_0$ forces $b$ to be a positive square and
$O>0$.  Indeed its two factors are coprime.  The congruence
$b\equiv1\pmod3$ excludes zero and negative squares.  A zero factor would
force $b=1$, hence $f=O=0$ and $1-4z^2=0$, which is impossible.
The remaining factor is a
square of the form $1+(3O-4)(bz)^2$ with $bz\ne0$, so the standard Pell
positivity lemma gives $O>0$.

For completeness, the positive bridge has $O\ge4$, so $d=3O-4>0$ is
nonsquare because it is 2 modulo 3.  A positive Pell solution
$u^2-dv^2=1$ can be chosen with $b\mid v$: raise a fundamental Pell unit
to the order of its image modulo $b$.  Taking $z=v/b>0$ makes $r_0$ a
square.  Sun's interpolation theorem then chooses $\xi$ divisible by
this $z$, so $x=\xi/z$ is still an integer.  No witness is added.

\paragraph{Smaller separating weights.}
For the actual polynomial $r_0$, $z\ne0$ implies $|r_0|>1$.  If $r_0=0$
or $r_0=\pm1$, its factorization forces $b=1$, hence $f=O=0$; but then
$r_0=1-4z^2\le-3$, a contradiction.  Also
\[
 \sqrt{|r_1|}\le P_n^2+K^2+2,
\]
because $|r_1|\le(P_n^2+4)K^2+4$ and the difference from the square of
the right-hand side is
$P_n^4+K^4+P_n^2K^2+4P_n^2\ge0$.
Thus the possibly signed weights
\[
 w_0=1,\qquad w_1=r_0,\qquad w_2=r_0(P_n^2+K^2+3)
\]
satisfy
\[
 |w_1|>\sqrt{|r_0|},\qquad
 |w_2|>\sqrt{|r_0|}+|w_1|\sqrt{|r_1|}.
\]
If a signed sum of these weighted roots is rational, any Galois
automorphism changing a root has a largest changed index.  Its term has
magnitude at least the corresponding weight and cannot be cancelled by
the earlier terms.  Hence every root is rational, and the integer
radicands are squares.  The same argument covers negative radicands by
complex absolute values, and zero radicands contribute no term.

\subsection{The final polynomial, without radicals or division}

Set
\[
 X_n=\Sigma m-\Theta,\qquad
 A_i=\Sigma^2w_i^2r_i\quad(0\le i\le2),
\]
and then
\[
 V_n=X_n^2+A_2-A_0-A_1,\qquad
 W_n=V_n^2+4X_n^2A_2-4A_0A_1.
\]
The requested equation is
\begin{equation}
 \boxed{F_{10}=W_n^2-16V_n^2X_n^2A_2=0.}
 \label{eq:teninteger}
\end{equation}
All preceding letters are polynomial abbreviations, not additional
unknowns.  Its formal norm identity is
\[
 F_{10}=\prod_{\varepsilon\in\{\pm1\}^3}
 \left(X_n+\sum_{i=0}^2\varepsilon_i\sqrt{A_i}\right).
\]
This identity explains the construction; equation~\eqref{eq:teninteger}
itself contains no radicals.  Its 35-term formal expansion is included
in the archive and checked exactly against the eight signed factors.

For a nonnegative zero of $P_0$, the positive bridge, the merged-square
Pell construction and the full residue-interpolation theorem supply
$r_0,r_1,r_2$ square, $\Sigma\ne0$ and $\Sigma\mid\Theta$.  Choose the
integer $m$ to make a norm factor vanish.  Conversely, if $F_{10}=0$,
then $\Sigma=0$ would give $F_{10}=\Theta^8\ne0$.  Thus $\xi=zx\ne0$
and $z\ne0$, so the structural separating-weight argument applies.
A vanishing norm factor forces the three radicands to be squares and
$\Sigma\mid\Theta$.  The quintic merger gives $S\mid T$ and the
coprime square gives $b$ square and $O>0$.  Sun's reverse bridge gives
a nonnegative zero of $P_0$.

Combining this with the fixed parent/compiler endpoint yields
\[
 \mathrm{ZFC}\vdash
 \left[\neg\exists(f,g,h,k,l,w,x,y,z,m)\in\mathbb Z^{10}\ F_{10}=0\right]
 \longleftrightarrow\operatorname{Con}(\mathrm{ZFC}).
\]
This invokes the stated mathematical theorems, including Sun's 2026
preprint.  The computational receipt verifies finite coefficient and
circuit claims; it is not an end-to-end formal proof of this equivalence.

\subsection{Reproduction and quantitative comparison}

The two commands \texttt{python build\_signed10.py} and
\texttt{python verify\_signed10.py}, run in the \texttt{signed10} directory,
use only the Python standard library.  The source sum is evaluated by a
shared-power Horner circuit; the geometric sum uses repeated doubling,
without division.  All constant powers are expanded into binary gates.
The main circuit has 566 additions, 25 subtractions and 851 multiplications,
with 79 fixed leaves.  A nonzero leading coefficient after diagonal
specialization modulo 101, 1009 and 1000003 proves that the circuit degree
upper bound is attained.

For transparency, the main degree calculation is as follows.  Here
$t=5^{31}$, and $X_n,A_i$ refer to the finalizer rather than the source
polynomial $X$:
\begin{center}
\begin{tabular}{@{}lr@{}}
\toprule
Expression & Total degree\\
\midrule
$X$ & $120t+24$\\
$Y$ & $80t+16$\\
$S$ & $280t+59$\\
$r_0$ & $1132t+250$\\
$r_1=R_2$ & $1920t+402$\\
$r_2=DFI^*$ & $24640t+5268$\\
$\Sigma$ & $1400t+297$\\
$X_n$ & $1400t+310$\\
$A_0$ & $3932t+844$\\
$A_1$ & $6984t+1496$\\
$A_2$ & $31624t+6766$\\
\bottomrule
\end{tabular}
\end{center}
The degree of $A_2$ strictly exceeds those of $X_n^2,A_0,A_1$.
Consequently the top part of the final norm is $A_2^4$, giving
$\deg F_{10}=126496t+27064$.  The recorded nonzero diagonal coefficients
are exact finite certificates that these top parts do not cancel for
the actual fixed coefficient tuple.

\begin{center}
\begin{tabular}{lrr}
Variant & Operations & Exact degree, $t=5^{31}$\\ \hline
Sun's four-square sextic specialization &1492&$2407680t+514640$\\
Four squares, quintic merger &1492&$2403200t+513696$\\
Three squares, Sun's generic weights &1445&$807360t+172560$\\
Three squares, universal successive weights &1444&$158592t+33880$\\
Three squares, structural signed weights &1442&$126496t+27064$
\end{tabular}
\end{center}
All five variants have the same ten unrestricted integer unknowns and
the same fixed source coefficients.  The complete proof, exact circuits,
coefficient recipes, binary payloads, source references and verification
receipts are in the \texttt{signed10} directory.

\section{A sixteen-integer equation with a smaller arithmetic circuit}
\label{sec:signed16}

This section gives an explicit intermediate tradeoff between the direct
Pell/compiler equations and the smaller-variable construction elsewhere in
this report.  The equation has sixteen unrestricted integer unknowns,
degree
\[
  813056\cdot5^{32}+10016
  =18930435180664062500000010016,
\]
and an evaluator with \(1617\) binary arithmetic operations.  It combines
package 30's 31-variable quartic, package 29's signed Jones interface and
relation-combining construction, and package 27's digit-polynomial
evaluation identities \cite{repo27,repo29,repo30}.  The new steps are a cofinal restriction of the
coding base, the removal of its sign guard, smaller root-separation weights,
and an automatic nonzero condition.  No minimum-variable or global
operation-minimality claim is made.

Every abbreviation below denotes a polynomial expression, not another
quantified unknown.  Every exponent in the equation is a fixed natural
number.  The large fixed coefficients are defined by finite recipes; they
are not existential parameters.

\subsection{The fixed instance and the upper-bound coefficient}

Let \(Q_{31}\) be the sum of the 22 squared quadratic residuals in the
\emph{unshifted-transport} parent
\texttt{signed16/data/parent\_quartic31.json}.  This is the package-30 parent,
with its whole four-port compiler tuple instantiated by the
\(121\)-by-\(8801\) construction of this report.  In particular, this parent
is not silently replaced by the locally shifted quartic evaluator.
The coefficient interface is the same polynomial interface in
\((\beta,K_0,M_C,M_F)\); all four ports use this one complete tuple.

Put \(\nu=31\), \(L=5^{32}\), and define the normalized polynomial
\[
 {\cal P}(x,n_1,\ldots,n_\nu)
       =Q_{31}(n_1+1,\ldots,n_\nu+1).
\]
The dummy input \(x\) is unused.  The degree is four, and
\({\cal P}(x,0,\ldots,0)>0\): the fifth parent residual is \(1\) when all
positive parent coordinates equal \(1\).
Write
\begin{align*}
 {\cal P}(x,n_1,\ldots,n_\nu)
    &=\sum_{\substack{\alpha\in\N^{\nu+1}\\|\alpha|\leq4}}
       p_\alpha (x,n_1,\ldots,n_\nu)^\alpha,\\
 c_\alpha&=(4-|\alpha|)!\left(\prod_{j=0}^{\nu}\alpha_j!\right)p_\alpha,
 \qquad
 \operatorname{exp}(\alpha)=\sum_{j=0}^{\nu}\alpha_j5^j .
\end{align*}
There are 402 nonzero symbolic weighted-coefficient records.  The omitted
records mean \(c_\alpha=0\), not omitted baseline digits.  The complete
finite coefficient bounds supplied with the new tuple establish
\[
 |c_\alpha|+1<2^{4269855}\qquad (|\alpha|\leq4).
\]
Set
\begin{align}
 z_*&=2^{4269855},&
 e_*(T)&=\sum_{|\alpha|\leq4}
        (z_*+c_\alpha)T^{L-\operatorname{exp}(\alpha)},\notag\\
 \upsilon_*&=\sum_{j=1}^{\nu}(2z_*)^{5^j},&
 y_*&=e_*(2z_*),&
 \kappa&=(2z_*)^{L+1}.
 \label{eq:s16-index}
\end{align}
Thus, exactly,
\[
 \kappa
 =2^{4269856(5^{32}+1)}
 =2^{99415332078933715820316769856}.
\]
The final polynomial uses \(z_*\), \(\kappa\), and the weighted coefficients
as fixed leaves.  Neither \(y_*\) nor \(\upsilon_*\) is a variable or a leaf
needed by its evaluator.

\begin{lemma}[The fixed upper bound]\label{lem:s16-kappa}
The recipes in \eqref{eq:s16-index} define positive integers and
\(y_*<\kappa\).  They give the Jones index for the normalized parent
\({\cal P}\) with parameter \(\nu=31\).
\end{lemma}
\begin{proof}
Each digit \(z_*+c_\alpha\) lies strictly between \(1\) and \(2z_*\).
Since every \(\alpha_j\leq4\), uniqueness of base-five expansions makes
the positions \(\operatorname{exp}(\alpha)\) distinct.  Also
\(\operatorname{exp}(\alpha)\leq4\cdot5^\nu<L\).
Consequently \(y_*\) is a base-\(2z_*\) integer whose positions lie in
\(\{0,\ldots,L\}\), with each digit less than \(2z_*\).  Therefore
\[
 0<y_*
 \leq(2z_*-1)\sum_{r=0}^{L}(2z_*)^r
 =(2z_*)^{L+1}-1<\kappa.
\]
The coefficient inequality, the power-of-two choice of \(z_*\), and the
displayed definitions of \(\upsilon_*,y_*\) are exactly the admissible
index conditions used by the inherited parametric Jones theorem.  This
uses that theorem at \(\nu=31\), not the convenience existence theorem
whose internal parent has 58 variables.
\end{proof}

\subsection{The literal sixteen-coordinate equation}

The unrestricted integer coordinates, in the emitted order, are
\begin{equation}
 t,g_0,h_0,s_0,w_0,\phi,i_0,j_0,\quad
 z_1,z_2,z_3,z_4,z_5,z_6,\quad v,u .
 \label{eq:s16-coordinates}
\end{equation}
First define
\[
 b=\kappa t^2,\qquad
 g=g_0+1,\quad h=h_0+1,\quad s=s_0+1,\quad
 w=w_0+1,\quad i=i_0+1,\quad j=j_0+1,
\]
and the fixed-degree coding polynomials
\begin{gather}
 {\cal B}=b^5,\quad q={\cal B}^{L},\quad
 \ell=\sum_{r=1}^{\nu}{\cal B}^{5^r},\quad
 \lambda=\sum_{r=0}^{4L-1}{\cal B}^{r},\quad
 e=e_*({\cal B}),\quad \theta={\cal B}-2z_*,\notag\\
 {\cal N}=q^{16},\qquad G_{\rm cap}=32b{\cal B}^{5^\nu}.
 \label{eq:s16-coding}
\end{gather}
The packed rank is
\begin{align}
 R={}&
 \left[
 g+eq^3+\ell q^5+
 \left\{
  2(e-z_*\lambda)(1+{\cal B}+g)^4+
  \lambda{\cal B}+\lambda{\cal B}q^4
 \right\}q^7
 \right]({\cal N}^2-{\cal N})\notag\\
 &+
 \left[
 q^3-b\ell+\ell+\theta\lambda q^3+({\cal B}-2)q^8
 \right]({\cal N}^2-1).
 \label{eq:s16-rank}
\end{align}
Continue with the unchanged Jones ratio expressions:
\begin{align}
 U&={\cal N}^2w,& Y&={\cal N}^2s,& M&=RY,&
 P&=2M^2U,\notag\\
 K&=R+1+h(P-1),&
 A&=M(U+1),&
 B_{\rm J}&=2R+1,& W_{\rm J}&=bw,\notag\\
 C&=B_{\rm J}+W_{\rm J}+\phi,&
 J&=4A-5,&
 Q&=3W_{\rm J}C-2(W_{\rm J}^2-1),\notag\\
 D&=(A^2-1)C^2+1,&
 E&=iJDC^2,&
 F&=(A^2-1)E^2+1,\notag\\
 G_1&=A+F(CD+1-A),&
 H&=B_{\rm J}+jC,&
 I&=(G_1^2-1)H^2+1.
 \label{eq:s16-ratio-block}
\end{align}
Finally put
\begin{align}
 a_0&=(P^2-1)K^2+1,&
 a_1&=DFI,\notag\\
 m&=K^2G_{\rm cap}
       -4(C-KY)^2G_{\rm cap}-8gK^2,&
 D_0&=JF,&
 T_0&=QF+J(H-C).
 \label{eq:s16-tests}
\end{align}
The auxiliary letter \(m\) denotes the margin in this section.

The five base coordinates that retain explicit sign guards are
\[
 (x_1,x_2,x_3,x_4,x_5)=(g_0,s_0,w_0,\phi,j_0).
\]
Starting with \(C_0=1\), form the six-stage stack
\begin{align}
 d_r&=3x_r+2,&
 B_r&=1+4d_rC_{r-1}^2z_r^2,&
 C_r&=C_{r-1}B_r &&(1\leq r\leq5),\notag\\
 d_6&=3m-1,&
 B_6&=1+4d_6C_5^2z_6^2,&
 C_*&=C_6=C_5B_6 .
 \label{eq:s16-stack}
\end{align}
The circuit names the stack outputs \texttt{guard1}, \ldots,
\texttt{guard6}; the ratio-block expression \(B_{\rm J}=2R+1\)
is separate.

For arbitrary integer arguments define
\begin{equation}
 {\cal R}(D',T',v)=
 \big((T'+D'v)^3-D'^2(T'+D'v)-D'^3\big)
 \big((T'+D'v)^2+23D'^2\big).
 \label{eq:s16-homogeneous}
\end{equation}
Using this polynomial, set
\begin{align}
 Z&=i\prod_{r=1}^{6}z_r,&
 S&=ZD_0^5,& T&={\cal R}(D_0,T_0,v),\notag\\
 w_1&=C_*,&
 V_0&=P^2+K^2+2,&
 w_2&=C_*V_0,\notag\\
 X&=Su-T,&
 A_0&=S^2C_*,&
 A_1&=S^2w_1^2a_0,&
 A_2&=S^2w_2^2a_1,\notag\\
 V&=X^2+A_2-A_0-A_1,&
 W&=V^2+4X^2A_2-4A_0A_1 .
 \label{eq:s16-finalizer}
\end{align}
The promised ordinary integer polynomial is
\begin{equation}
 \boxed{\displaystyle F_{16}=W^2-16V^2X^2A_2.}
 \label{eq:s16-polynomial}
\end{equation}
In particular, \(t\) is deliberately absent from \(Z\).  Its nonvanishing
will follow from the recovered margin; omitting it saves one multiplication.

\subsection{The exact imported interface and the new cofinal restriction}

For an integer \(a\), write
\(\operatorname{Sq}(a)\) for the assertion
\(\exists r\in\Z\;(a=r^2)\).
The inherited Jones ratio equivalence applies to the eight
\emph{nonnegative} coordinates
\[
 \varepsilon,g_0,h_0,s_0,w_0,\phi,i_0,j_0,
 \qquad b=y_*+1+\varepsilon.
\]
For the fixed normalized parent, it says that
\begin{equation}
 \operatorname{Sq}(a_0),\quad \operatorname{Sq}(a_1),\quad
 D_0\mid T_0,\quad m>0
 \label{eq:s16-jr}
\end{equation}
are jointly solvable if and only if \({\cal P}\) has a nonnegative
31-tuple zero.  The ratio expressions in
\eqref{eq:s16-coding}--\eqref{eq:s16-tests} are precisely those in this
interface.  This section does not apply a signed interpretation to an
unspecified nine-variable polynomial.

The following strengthened completeness fact is needed to replace
\(b=y_*+1+\varepsilon\) by \(b=\kappa t^2\).

\begin{lemma}[Cofinal powers of two]\label{lem:s16-cofinal}
Given a nonnegative parent witness
\({\cal P}(1,n_1,\ldots,n_\nu)=0\), the Jones ratio conditions
\eqref{eq:s16-jr} have witnesses with \(b=2^a\) for every sufficiently
large integer \(a\).  In particular, they have witnesses with
\(b=\kappa t^2\) and \(t>0\).
\end{lemma}
\begin{proof}
Let \(\xi=(1,n_1,\ldots,n_\nu)\).
The constructive master-system proof
\texttt{Master3.USys\_of\_mem} chooses
\[
 a_0'=y_*+\sum_{r=0}^{\nu}\xi_r+1,\qquad b=2^{a_0'}.
\]
Inspection of that complete proof shows that the equality specifying
\(a_0'\) is used only to obtain
\[
 a\geq1,\qquad y_*<b,\qquad \xi_r<b\quad(0\leq r\leq\nu).
 \tag{C}
\]
These conclusions hold for every \(a\geq a_0'\), because
\(a<2^a\).  The rest of the construction depends on (C), on \(b=2^a\),
and on the given polynomial zero; it imposes no upper bound on \(a\).

More explicitly, at this chosen base one takes the polynomial evaluations
\[
 e=e_*({\cal B}),\qquad
 \ell=\sum_{r=1}^{\nu}{\cal B}^{5^r},\qquad
 g=\sum_{r=1}^{\nu}\xi_r{\cal B}^{5^r},\qquad {\cal B}=b^5,
\]
together with \(q,\lambda,\theta\) from \eqref{eq:s16-coding}.
Normalization of the parent implies \(g>0\).
The digit bounds and the source's carry-free coding identity give
\[
 e<2z_*{\cal B}^{L},\qquad
 \ell<2{\cal B}^{5^\nu},\qquad
 g<b{\cal B}^{5^\nu}.
\]
The necessary positive slack is
\[
 \alpha_{\rm master}
       =(b-y_*)q^2-e\ell g^2.
\]
It is positive uniformly in the larger base: putting \(r=5^\nu\),
\[
 e\ell g^2
 <4z_*b^2{\cal B}^{L+3r}
 \leq{\cal B}^{L+3r+1}
 \leq{\cal B}^{2L}=q^2.
\]
Here \(4z_*b^2\leq b^5={\cal B}\) follows from \(b>y_*\geq2z_*\),
and \(L+3r+1\leq2L\) follows from \(L=5r\).
The two master quotients
\[
 \frac{e-y_*}{\theta},\qquad
 \frac{\ell-\upsilon_*}{\theta}
\]
are positive integers: divisibility is polynomial evaluation modulo
\({\cal B}-2z_*\), and positivity follows from
\({\cal B}>2z_*\).  The three bit-disjointness identities in the master
system are the inherited Jones coding identities; their hypotheses are
(C), \(b=2^a\), and the unchanged parent zero.  Thus the same explicit
construction works at every such larger base.

The proof of \texttt{NineReduce.reduced\_of\_mem} keeps this very base
and gives the required packed-rank/binomial condition, as well as
\(g<b{\cal B}^{5^\nu}\).  The theorem
\texttt{NineRatio.ratio\_complete} then takes \(b\) as a fixed parameter;
its hypothesis \(32g<G_{\rm cap}\) follows from that last inequality.
It supplies the positive \(h,s,w,i,j\) and nonnegative \(\phi\) in
\eqref{eq:s16-jr}, without selecting a different base.

Write \(\kappa=2^c\), where
\(c=4269856(L+1)\).  Choose any \(r\geq0\) with
\(c+2r\geq a_0'\) and set \(t=2^r\).  Then
\[
 \kappa t^2=2^{c+2r}
\]
is one of the permitted bases.  This gives the required positive \(t\)
and proves the assertion constructively.
\end{proof}

\subsection{Signed coordinates whose guards are unnecessary}

\begin{lemma}[A denominator that never vanishes]\label{lem:s16-denominator}
For every integer assignment of the eight base coordinates,
\(D_0=JF\neq0\).
\end{lemma}
\begin{proof}
The integer \(J=4A-5\) is nonzero.  If
\(F=(A^2-1)E^2+1=0\), then
\((1-A^2)E^2=1\), so \(A=0\) and \(E=\pm1\).
But \(J=-5\) when \(A=0\), and \(E=iJDC^2\) is divisible by \(5\).
This is impossible.
\end{proof}

\begin{lemma}[The sign of \(i\)]\label{lem:s16-isign}
If \(i\neq0\), replacing \(i\) by \(|i|\) preserves
\(a_0,a_1,D_0,T_0,m\), and makes \(i_0=|i|-1\) nonnegative.
\end{lemma}
\begin{proof}
The only unsquared use of \(i\) is in \(E=iJDC^2\).
Among the five expressions under consideration, every dependence on
\(E\) passes through \(E^2\) in \(F\).  Thus they depend only on \(i^2\).
\end{proof}

\begin{lemma}[The \(h\) guard follows from the margin]\label{lem:s16-hguard}
Suppose \(b>y_*\), the five coordinates
\(g_0,s_0,w_0,\phi,j_0\) are nonnegative, \(a_0\) is an integer square,
and \(m>0\).  Then \(h\geq1\), hence \(h_0\geq0\).
\end{lemma}
\begin{proof}
We have \(g,s,w,j\geq1\), \(\phi\geq0\),
\(b\geq2\), \({\cal B}=b^5>2z_*\), \(e\geq0\),
and \(G_{\rm cap}>0\).
The margin inequality first gives \(K\neq0\) and
\begin{equation}
 g<G_{\rm cap}/8=4b{\cal B}^{5^\nu},
 \qquad
 4(C-KY)^2<K^2.
 \label{eq:s16-margin-consequences}
\end{equation}
Indeed \(m\leq K^2(G_{\rm cap}-8g)\); then discarding the strictly
positive term \(8gK^2\) gives the second inequality.

Put \(r=5^\nu\).  Since \({\cal B}>5\) and \(b\leq{\cal B}\),
\[
 1+{\cal B}+g
 <1+{\cal B}+4{\cal B}^{r+1}
 \leq5{\cal B}^{r+1}
 <{\cal B}^{r+2}
 \leq{\cal B}^{5r}=q.
\]
Also \(\ell<{\cal B}^{r+1}\), hence
\(b\ell<{\cal B}^{r+2}\leq q<q^3\).
The potentially negative part of the first bracket in
\eqref{eq:s16-rank} is absorbed by its positive term:
\[
 \lambda{\cal B}q^4
       >2z_*\lambda(1+{\cal B}+g)^4 .
\]
All the remaining terms of that bracket are nonnegative, and \(g>0\).
The second bracket is positive because \(q^3>b\ell\), \(\theta>0\),
and \({\cal B}>2\).
It follows that
\[
 R\geq{\cal N}^2-{\cal N}\geq2.
\]
Thus \(U,Y\geq1\), \(C=2R+1+bw+\phi>0\), and
\[
 P=2R^2Y^2U\geq2R^2>R+2.
\]

The second margin consequence rules out \(K\leq0\):
\(K=0\) is impossible, and \(K<0\) would give
\(|C-KY|=C+|K|Y>|K|\).
Therefore \(K>0\).
If \(h\leq0\), the relation \(K=R+1+h(P-1)\) gives
\(1\leq K\leq R+1<P\).
For \(K=1\) it would imply \(P-1\mid R\), contradicting
\(0<R<P-1\).
For \(2\leq K<P\),
\[
 (PK-1)^2<(P^2-1)K^2+1<(PK)^2,
\]
since the two differences are \(K(2P-K)>0\) and \(K^2-1>0\).
This contradicts squarehood of \(a_0\).
Consequently \(h\geq1\).
\end{proof}

\subsection{The coprime Pell stack}

\begin{lemma}[Six guards in one square test]\label{lem:s16-stack}
Subject to \(z_1\cdots z_6\neq0\), squarehood of \(C_*\) in
\eqref{eq:s16-stack} implies
\begin{equation}
 g_0,s_0,w_0,\phi,j_0\geq0,\qquad m>0.
 \label{eq:s16-guards}
\end{equation}
Conversely, if these inequalities hold, the six \(z_r\) can be chosen
nonzero so that \(C_*\) is a square.
\end{lemma}
\begin{proof}
Every \(B_r\equiv1\pmod4\), so no \(B_r\) vanishes.
Also \(\gcd(B_r,C_{r-1})=1\).
Inductively, the six factors \(B_r\) are pairwise coprime in absolute
value.  If their product is a square, unique prime factorization shows
that every \(|B_r|\) is a square.  A negative square is \(0\) or \(3\)
modulo \(4\), so \(B_r\equiv1\pmod4\) forces every factor to be a
positive square.  If any \(d_r<0\), the nonzero integers
\(C_{r-1},z_r\) would give
\(B_r=1+4d_rC_{r-1}^2z_r^2\leq-3\), a contradiction.
Finally \(d_r\equiv2\pmod3\), so it cannot be zero.
Thus \(d_r>0\) for all \(r\).  For \(r\leq5\) this means \(x_r\geq0\);
for \(r=6\) it means \(m\geq1\).

For the converse, each \(d_r>0\) and \(d_r\equiv2\pmod3\), so it is
not a square.  The Pell equation
\[
 a^2-d_r v^2=1
\]
has a nontrivial positive solution.  It has a positive solution with
\(v\) divisible by any prescribed positive integer \(M\): the matrix
\(\left(\begin{smallmatrix}a&d_rv\\v&a\end{smallmatrix}\right)\)
has determinant \(1\), and some positive power of its reduction modulo
\(M\) is the identity in the finite group
\(\mathrm{SL}_2(\Z/M\Z)\).
Choose \(M=2|C_{r-1}|\), and put \(z_r=v/(2C_{r-1})\).
Then \(z_r\) is a nonzero integer and \(B_r=a^2\).
Starting with \(C_0=1\), induction makes every \(C_r\) a positive square.
\end{proof}

A useful strengthening is available before any guard inequality is known.
If all \(z_r\neq0\), then \(d_r\neq0\), \(C_{r-1}\neq0\), and hence
\(B_r\neq1\).  The congruence \(B_r\equiv1\pmod4\) also excludes
\(B_r=-1\).  Therefore
\begin{equation}
 |C_*|>1
 \quad\hbox{whenever}\quad z_1\cdots z_6\neq0.
 \label{eq:s16-stack-large}
\end{equation}
This preguard estimate is what permits the smaller weights.

\subsection{Divisibility and nonvanishing in one relation}

\begin{lemma}[An intersective polynomial]\label{lem:s16-intersective}
The polynomial
\[
 \Pi(T)=(T^3-T-1)(T^2+23)
\]
has no rational zero, but has an integer zero modulo every positive
integer modulus.
\end{lemma}
\begin{proof}
The only possible rational roots of the monic cubic are \(1,-1\),
and neither is a root.  The quadratic has no real root.

For an odd prime \(p\neq23\), the cubic has discriminant \(-23\).
If it has a root modulo \(p\), the root is simple and lifts to every
power of \(p\).
If it has no such root, it is irreducible of degree three.
Frobenius cycles its three roots in \(\mathbb F_{p^3}\), an even
permutation, so their Vandermonde product lies in \(\mathbb F_p\).
Its square is \(-23\), which is nonzero.  Thus \(T^2+23\) has a
simple root modulo \(p\), again lifting to every power.
At \(p=23\), the cubic has the simple root \(3\), since
\(3^3-3-1=23\) and \(3(3^2)-1=26\not\equiv0\pmod{23}\).

For \(p=2\), use \(-23\equiv1\pmod8\).
If an odd \(a\) is a root of \(T^2+23\) modulo \(2^e\), \(e\geq3\),
one of \(a\) and \(a+2^{e-1}\) is a root modulo \(2^{e+1}\):
their squares differ by \(2^ea\) modulo \(2^{e+1}\).
This supplies roots for all powers of two.
The Chinese remainder theorem combines roots for the prime-power
factors of any positive modulus.
\end{proof}

\begin{lemma}[Homogeneous divisibility merger]\label{lem:s16-merger}
For integers \(D',T',Z'\) with \(D'\neq0\),
\begin{equation}
 \big(D'\mid T'\ \hbox{and}\ Z'\neq0\big)
 \quad\Longleftrightarrow\quad
 \exists v\in\Z:\ Z'D'^5\mid{\cal R}(D',T',v).
 \label{eq:s16-merger}
\end{equation}
Moreover \({\cal R}(D',T',v)\neq0\) for every integer \(v\).
\end{lemma}
\begin{proof}
The rational identity
\[
 {\cal R}(D',T',v)=D'^5\Pi(T'/D'+v)
\]
proves nonvanishing by Lemma~\ref{lem:s16-intersective}.
If \(T'=D'q\) and \(Z'\neq0\), choose \(v\) so that
\(Z'\mid\Pi(q+v)\); this proves the forward implication.

Conversely, \(Z'=0\) would make the displayed divisibility assert
\({\cal R}=0\), which is impossible.
If \(D'\nmid T'\), there is a prime \(p\) with
\(\alpha=v_p(|D'|)>\beta=v_p(|T'|)\).
Here \(T'\neq0\).
For \(r=T'+D'v\), one has \(v_p(r)=\beta\).
In \(r^3-D'^2r-D'^3\), the first term has the unique least valuation
\(3\beta\); in \(r^2+23D'^2\), it has the unique least valuation
\(2\beta\).  Hence
\[
 v_p({\cal R})=5\beta<5\alpha=v_p(D'^5),
\]
contradicting the asserted divisibility.  Signs of \(D'\) and \(Z'\)
do not affect this argument.
\end{proof}

\subsection{The three-root norm with signed weights}

\begin{lemma}[The finalizer]\label{lem:s16-finalizer}
Fix all coordinates except \(u\).
For the expressions above, \(\exists u\in\Z:F_{16}=0\) is equivalent to
\begin{equation}
 S\neq0,\qquad S\mid T,\qquad
 \operatorname{Sq}(C_*),\quad
 \operatorname{Sq}(a_0),\quad
 \operatorname{Sq}(a_1).
 \label{eq:s16-finalizer-interface}
\end{equation}
\end{lemma}
\begin{proof}
Lemma~\ref{lem:s16-denominator} gives \(D_0\neq0\), and
Lemma~\ref{lem:s16-merger} gives \(T\neq0\).
If \(S=0\), then \(A_0=A_1=A_2=0\), \(X=-T\), and
\(F_{16}=T^8\neq0\).
Thus every zero has \(S\neq0\).
Since \(S=i(z_1\cdots z_6)D_0^5\), every \(z_r\) and \(i\) is
nonzero; in particular \eqref{eq:s16-stack-large} applies.

For all real \(P,K\),
\begin{equation}
 \sqrt{|(P^2-1)K^2+1|}\leq P^2+K^2+1 .
 \label{eq:s16-a0-bound}
\end{equation}
Indeed
\(|(P^2-1)K^2+1|\leq(P^2+1)K^2+1\), and the difference between
\((P^2+K^2+1)^2\) and this last upper bound is
\[
 P^4+K^4+P^2K^2+2P^2+K^2\geq0.
\]
It follows from \(|C_*|>1\) that the possibly signed weights satisfy
\begin{align}
 |w_1|&>\sqrt{|C_*|},\notag\\
 |w_2|&>\sqrt{|C_*|}+|w_1|\sqrt{|a_0|}.
 \label{eq:s16-separation}
\end{align}
These inequalities require neither positive radicands nor recovered
guard inequalities.

Choose complex roots
\(\rho_0^2=C_*,\rho_1^2=a_0,\rho_2^2=a_1\).
Algebraic multiplication gives
\begin{equation}
 F_{16}
  =\prod_{\epsilon_0,\epsilon_1,\epsilon_2\in\{-1,1\}}
     \left[
       X-S(\epsilon_0\rho_0+
             \epsilon_1w_1\rho_1+\epsilon_2w_2\rho_2)
     \right].
 \label{eq:s16-eight-factor}
\end{equation}
For example, first multiplying the \(\rho_0,\rho_1\) signs gives
\(\big((X\mp\sqrt{A_2})^2-A_0-A_1\big)^2-4A_0A_1\);
the two remaining factors are
\(W\mp4VX\sqrt{A_2}\), whose product is
\eqref{eq:s16-polynomial}.

A zero makes one of the signed sums in
\eqref{eq:s16-eight-factor} equal to \(X/S=u-T/S\), a rational number.
Distinct sign patterns on nonzero roots give distinct sums:
at the largest index where signs differ, the absolute value of the
weighted root is at least the absolute value of its weight, since a
nonzero integer radicand has root magnitude at least \(1\).
Inequalities \eqref{eq:s16-separation} make this strictly larger than
the sum of the magnitudes of all preceding weighted roots.
The triangle inequality therefore precludes cancellation.  This remains
valid for imaginary roots of negative radicands and for negative weights.
Zero roots may simply be omitted from this comparison.

The field \(\mathbb Q(\rho_0,\rho_1,\rho_2)\) is the splitting field of
three quadratic polynomials, hence a finite Galois extension of
\(\mathbb Q\).  Each automorphism changes root signs.
Since the encoded sum is rational, every automorphism fixes it.
The separation just proved forces every nonzero root in it to be fixed
individually.  Therefore all three roots are rational.
A rational root of a monic polynomial \(T^2-a\), \(a\in\Z\), is an
integer.  Thus \(C_*,a_0,a_1\) are nonnegative integer squares.
The encoded sum is now an integer, so \(T/S\) is an integer as well.
This proves necessity of \eqref{eq:s16-finalizer-interface}.

Conversely, under that interface choose integer square roots and set
\[
 u=T/S+\rho_0+w_1\rho_1+w_2\rho_2 .
\]
This is an integer and makes a factor in
\eqref{eq:s16-eight-factor} vanish.
\end{proof}

\subsection{Soundness, completeness, and the consistency equivalence}

\begin{theorem}[The sixteen-integer presentation]\label{thm:s16}
For the fixed coefficients above, ordinary mathematics formalizable
in ZFC proves
\[
 \big(\exists {\bf x}\in\Z^{16}:F_{16}({\bf x})=0\big)
 \ \longleftrightarrow\
 \big(\text{there is a finite syntactic ZFC refutation}\big).
\]
Consequently,
\[
 \mathrm{ZFC}\ \vdash\
 \big(\nexists{\bf x}\in\Z^{16}:F_{16}({\bf x})=0\big)
 \ \longleftrightarrow\ \operatorname{Con}(\mathrm{ZFC}).
\]
\end{theorem}
\begin{proof}
For soundness, let \(F_{16}=0\).
Lemma~\ref{lem:s16-finalizer} gives \(S\neq0\), \(S\mid T\), and
the three square conditions.
Lemma~\ref{lem:s16-merger} gives \(D_0\mid T_0\) and
\(i z_1\cdots z_6\neq0\).
Lemma~\ref{lem:s16-stack} recovers the five nonnegativity guards
and \(m>0\).

At this point \(t\neq0\) follows without a separate nonzero factor.
If \(t=0\), then \(b=\kappa t^2=0\) and \(G_{\rm cap}=0\), while
\(g=g_0+1\geq1\).  The literal margin formula would give
\[
 m=-8gK^2\leq0,
\]
a contradiction.
Hence \(b=\kappa t^2\geq\kappa>y_*\), and
\(\varepsilon=b-y_*-1\) is a nonnegative integer.
Lemma~\ref{lem:s16-hguard} gives \(h_0\geq0\).
Lemma~\ref{lem:s16-isign} replaces \(i\) by \(|i|\), preserving all
four Jones conditions and making \(i_0\geq0\).
We have therefore recovered an admissible nonnegative tuple for the
unchanged Jones interface \eqref{eq:s16-jr}.
The inherited ratio soundness theorem gives a zero of the parent
quartic.  Its complete native/Pell compiler interface gives a blank
K215 halting computation, and the calibrated K215 theorem gives a
finite syntactic ZFC refutation.

For completeness, a finite ZFC refutation gives a blank K215 halting
run and then a positive parent-quartic solution by those same two
interfaces.
After the coordinate shift to the normalized parent, apply
Lemma~\ref{lem:s16-cofinal}.
It gives a Jones tuple satisfying \eqref{eq:s16-jr} at
\(b=\kappa t^2\) with \(t>0\), all the ordinary base coordinates
nonnegative, and \(i>0\).
Lemma~\ref{lem:s16-stack} supplies nonzero \(z_1,\ldots,z_6\)
for which \(C_*\) is a square.
Now \(Z=i z_1\cdots z_6\neq0\); the inherited Jones tuple has
\(D_0\mid T_0\).
Lemma~\ref{lem:s16-merger} chooses \(v\) with \(S\mid T\).
The three radicands are squares, so
Lemma~\ref{lem:s16-finalizer} chooses the final integer \(u\)
and gives \(F_{16}=0\).

Every extra root, quotient, or exponent selected in this proof is a
proof witness used to establish equivalence.  The displayed equation
itself quantifies only the sixteen coordinates in
\eqref{eq:s16-coordinates}.
Negating the existence equivalence gives the stated consistency
equivalence.  Neither direction assumes or proves
\(\operatorname{Con}(\mathrm{ZFC})\).
\end{proof}

\subsection{Degree and arithmetic size}

All degrees in this subsection give each of the sixteen unknowns
degree one and every fixed coefficient degree zero.
The largest term in the packed rank is
\(\lambda{\cal B}q^{11}{\cal N}^2\);
it gives \(\deg R=470L\).
The largest term in the margin is
\(-4K^2Y^2G_{\rm cap}\), giving \(\deg m=4442L+12\).
The useful intermediate degrees are
\[
\begin{array}{c|r@{\qquad}c|r}
 R&470L&P&1900L+3\\
 K&1900L+4&A&1110L+2\\
 a_0&7600L+14&a_1&49280L+68\\
 m&4442L+12&C_5&363\\
 C_*&4442L+1103&D_0&13750L+20\\
 S&68750L+107&T&68750L+105\\
 X&68750L+108&w_2&8242L+1111\\
 A_0&141942L+1317&A_1&153984L+2434\\
 A_2&203264L+2504&&
\end{array}
\]
For the first five stack stages the degree recurrence is
\(c_r=3c_{r-1}+3\), starting at \(c_0=0\).
The last stage has degree
\(3\cdot363+(4442L+12)+2=4442L+1103\).
The seven factors of \(Z=i\prod_{r=1}^6z_r\) account for the constant
\(107=5\cdot20+7\) in the degree of \(S\).

The term \(A_2\) has strictly greater degree than \(X^2,A_0,A_1\).
Thus \(V\) has the leading form of \(A_2\), \(W\) that of \(A_2^2\),
and the final polynomial that of \(A_2^4\).
All of these leading forms are nonzero: the rank has the displayed
unique top term, the margin has its unique negative top term, and
subsequent relevant leading forms are products and squares of
nonzero forms.  In particular, after the authentic positive fixed
coefficients are substituted,
\[
 \deg F_{16}=4(203264L+2504)=813056L+10016.
\]
This is an exact degree assertion, not just a structural upper bound
that ignores cancellations.

The literal pruned evaluator contains
\[
 697\ \hbox{additions}
 \;+\;26\ \hbox{subtractions}
 \;+\;894\ \hbox{multiplications}
 \;=\;1617\ \hbox{operations}.
\]
Each binary operation, including a square and multiplication by a
fixed numeral, is charged.  Reuse and fixed leaves are free, consistently
with the repository's arithmetic-circuit convention.  Fixed exponentiation
is expanded into multiplication gates; no exponentiation oracle is used.

For clarity, the compact evaluation of the enormous digit baseline
does not omit its zero-coefficient positions.  If
\(p_r={\cal B}^{5^r}\), then the baseline is represented through the
complete homogeneous polynomial of degree four, and
\[
 \lambda
 =\left(\prod_{r=0}^{\nu}
       (1+p_r+p_r^2+p_r^3+p_r^4)\right)(1+q)(1+q^2).
\]
These are polynomial identities.
The full baseline contains
\(\binom{36}{4}=58905\) distinct positions; the 402 nonzero
weighted-coefficient records supply its sparse correction.

\subsection{Verification boundary}

The final binding is
\texttt{signed16/data/INSTANCE.json}; it names the whole compiler
tuple, parent, weighted-coefficient records, and index recipes.
The literal division-free polynomial is
\texttt{signed16/generated/signed16.json}.
The new reduction does not claim a fresh end-to-end Lean or Rocq
certificate.  The source Lean files explain the precise Jones interfaces
and the cofinal construction, while the native compiler and its K215
calibration remain identified mathematical proof dependencies.

The independent data-only auditor imports neither the emitter nor its DAG
utility.  It checks all four decompressed port hashes, enumerates the
58905 baseline positions directly, evaluates the geometric sum by a
separate method, and compares 288 named-node values on six signed
assignments over two primes, including \(t=0\).
It also checks the eight-factor norm in the rank-eight ring with three
formal square roots on 100 exact assignments.
For the final numerical instance, the top-degree diagonal coefficients
are nonzero modulo two primes:
\[
 121738\pmod{1000003},\qquad
 73024\pmod{1000033}.
\]
Together with the propagated degree upper bound, either nonzero residue
is an independent certificate that the top degree survives on this fixed
coefficient slice.

These checks establish transcription, coefficient binding, and algebraic
metrics.  The universal soundness and completeness claims rest on the
proofs above and the stated inherited interfaces, not on finite tests or
a search for a Diophantine solution.


\section{Smaller positive systems and circuit evaluations}
\label{sec:low-degree}

The coefficient adapter also supports two improvements that preserve the
positive-domain arithmetic constructions of package~30 \cite{repo30}.  One
retains degree four and reduces the arithmetic evaluation cost; the other
retains twelve unknowns and reduces the cost of its much higher-degree norm.
All the identities in this section hold with the four coefficient ports
treated as independent indeterminates.  We then specialize \emph{all four}
ports to the new complete $121$ by $8801$ tuple.  The resulting fixed
polynomials need not equal the historical fixed polynomials, because their
coefficients have changed.  Their arithmetic templates have the exact
identities proved below.

\begin{center}
\begin{tabular}{lrrrr}
\hline
Presentation & Unknowns & Equations & Degree & Operations\\
\hline
Quadratic comparison system &31&22&2&60\\
Its sum of squares &31&1&4&125\\
Twelve-variable product &12&1&1081&108\\
Twelve-variable system &12&3&1072&103\\
Transport-shifted system &12&3&1072&102\\
\hline
\end{tabular}
\end{center}

Every unknown in this table is a strictly positive integer.  For a system,
``degree'' means the maximum degree of a residual.  Every binary addition,
subtraction and multiplication is charged, including a multiplication by a
fixed numeral; a comparison and reuse of an existing register are free.
The source costs at the same two points were $65/130$ operations for the
quadratic system/quartic, and $107/111$ for the twelve-variable
system/product.  These are local improvements.  The separate inherited
$59$-operation higher-degree systems and $67$-operation single equation
retain their lower total-operation counts.

\subsection{A 60-operation quadratic system}

The $31$ supplied positive integers are
\[
\begin{gathered}
J,F,\alpha,z,f,h,i,o,j,s,w,\tau,\eta,\zeta,y,\gamma,\\
N,E,P,c,a,X,D,C_2,L,F_2,V,g,I,B_{\rm aux},\Delta.
\end{gathered}
\]
Use the fixed ports $\beta,K_0,M_C,M_F$, where $M_F$ is the already shifted
high mask.  Introduce only the polynomial abbreviations
\[
\begin{gathered}
q=\beta J+1,\quad k=\eta+\zeta,\quad H=4a+3,
\quad C=q-F-\alpha,\quad U=L+1,\\
Y=a-E,\qquad Q=c-\eta,\qquad T=P-k,\qquad R=jc-V.
\end{gathered}
\]
In particular, $R$ is quadratic.  The following are the $22$ comparisons:
\begin{equation}
\begin{aligned}
 U&=q^2,&N&=qU,&X&=wq,&Y&=sN,\\
 E&=XY,&Q&=kY,&T&=EQ,&\tau^2&=PT+1,\\
 D&=X+ac+\gamma H,&\Delta&=a^2+H,&C_2&=c^2,\\
 D^2&=\Delta C_2+1,&k&=R+1+hE.
\end{aligned}
\label{eq:quadratic-first}
\end{equation}
\begin{equation}
\begin{aligned}
 R&=(q-F)(N-q)-CL+(M_C+qM_F)J,\\
 (K_0+w)C&=F+z\beta J,\\
 F_2&=f^2,\qquad V=of-c,\qquad g=y^2-V^2,\\
 I&=iC_2,\qquad B_{\rm aux}=I^2,
       \qquad B_{\rm aux}=\Delta(F_2-1),\\
 y^2-B_{\rm aux}g&=1.
\end{aligned}
\label{eq:quadratic-second}
\end{equation}
Each equality has a quadratic residual after substituting the displayed
abbreviations.  If these residuals are $r_1,\ldots,r_{22}$ in the displayed
order, the single polynomial is
\begin{equation}
 \boxed{\qquad Q_{125}=\sum_{j=1}^{22}r_j^2.\qquad}
 \label{eq:q125}
\end{equation}
Thus its positive zero set is exactly the system's positive solution set.

The first four simplifications reorient existing equalities so that a
computed difference or sum serves two purposes.  The equality
$L=q^2-1$ becomes $L+1=q^2$, and the same $L+1$ occurs in $N=q(L+1)$.
The equality $a=E+sN$ becomes $a-E=sN$, reusing the $a-E$ in $E=X(a-E)$.
Similarly, $c-\eta=k(a-E)$ reuses $c-\eta$ in the equation defining $P$,
and $P-k=E(c-\eta)$ reuses $P-k$ in $\tau^2=P(P-k)+1$.
These are polynomial identities on all integer assignments.

The fifth simplification is the positive transport-coordinate shift also
used in package~26 \cite{repo26}.  Its original equation is
\[
 (K_0+w)C+(q-F)-z_{\rm old}(q-1)=1.
\]
Substituting $z_{\rm old}=z+1$ and $q-1=\beta J$ gives exactly the second
equation in \eqref{eq:quadratic-second}.  A new positive solution therefore
gives an old one.  Conversely, first apply the established compiler
reconstruction to an old solution.  It gives
\[
 C>0,\qquad F<q,\qquad X=wq=2^R,\qquad R>3q,
\]
and hence $w>q>F$.  The original transport equation implies
\[
 (z_{\rm old}-1)(q-1)=(K_0+w)C-F>0.
\]
Thus $z_{\rm old}\ge2$, and $z=z_{\rm old}-1$ remains positive.  This
establishes a bijection of positive zero sets.  The positivity claim is a
consequence of the completed source reconstruction; it is not presumed in
the initial bootstrap.

The resulting comparison circuit has $33$ multiplications and $27$
additions/subtractions, hence $60$ operations.  Forming the residuals,
squaring them and summing gives $55$ multiplications and $70$
additions/subtractions, hence $125$ operations.  Its degree is exactly
four: $\tau$ occurs only in the residual
$\tau^2-P(P-k)-1$, so the coefficient of $\tau^4$ in \eqref{eq:q125}
is exactly one.  This argument is uniform in every fixed coefficient
slice and does not depend on a generic numerical specialization.

The supplied builder checks exact equality of all $22$ residuals with
the source residuals after $z_{\rm old}=z+1$, treating all four ports
symbolically.  It expands the quartic to $299$ terms when port exponents
are included in the symbolic monomials.  Specializing actual coefficients
can identify terms, so $299$ is stated with that convention.  The
degree-four certificate is unaffected.

\subsection{A 108-operation twelve-variable equation}

The twelve supplied positive integers are now
\[
 J,F,\alpha,z,h,i,T,s,w,\eta,\zeta,t.
\]
The notation in this subsection is local.  Set
\begin{align*}
 q&=\beta J+1,& C_w&=q-F-\alpha,&X&=wq,&Y&=sq^3,&E&=XY,\\
 k&=\eta+\zeta,&c&=kY+\eta,&a&=E+Y,&H&=4a+3,&\Delta&=a^2+H,\\
 \ell&=EkY,&A&=1+\ell(\ell+k),&G&=1+\Delta c^2,&b_0&=X+ac.
\end{align*}
The transport and index registers are
\begin{align*}
 R&=((q-1)(q-F)+\alpha)(q^2-1)+(M_C+qM_F)J,\\
 N_k&=k-hE-R,\qquad
 N_t=(K_0+w)C_w+q-F-z(q-1).
\end{align*}
For the Pell norm define
\[
\begin{gathered}
 I=ic^2,\quad D=1+\Delta I^2,\quad Q=\Delta(D-1),\quad d=Q-1,
 \quad b=c+RD,\\
 M=k(I^2+T^2),\qquad L=HM+1,\qquad W=b_0+Ht.
\end{gathered}
\]
The economical scalar evaluator is
\begin{align*}
 S&=H^2Q,\qquad H_P=SDc^2T^2,\qquad H_U=Sb^2,\\
 n&=(HM)^2A,\qquad d_n=dn,\qquad H_e=H_P+H_U-H^2,\\
 j&=4H_PH_U,\qquad \mathcal T=2d(W^2+L^2G),\\
 r&=d(W^2-L^2G),\qquad\mathcal K=r^2,\\
 a_1&=\mathcal T-2(d_n+H_e),\qquad b_1=(d_n-H_e)^2-j-\mathcal K,\\
 v&=a_1(a_1\mathcal T+2b_1)-16d_nj,\\
 u&=a_1^2\mathcal K-b_1^2.
\end{align*}
All these letters denote polynomial abbreviations, not new witnesses.
Finally put
\begin{equation}
 N=\mathcal K v^2-u(\mathcal T v-u),\qquad
 \boxed{\quad P_{108}=N_kN_t(1+N)-1.\quad}
 \label{eq:p108}
\end{equation}
Its equation is $P_{108}=0$.  Its three-equation version is
\begin{equation}
 N_k=1,\qquad N_t=1,\qquad
 \mathcal K v^2=u(\mathcal T v-u).
 \label{eq:system103}
\end{equation}
Replacing the middle equation by
\begin{equation}
 (K_0+w)C_w=F+z(q-1)
 \label{eq:positive102transport}
\end{equation}
gives the $102$-operation system.  Its $z$ is the old transport coordinate
minus one; the positivity and equivalence proof is exactly the one above,
applied after the twelve-variable reconstruction.

\subsection{Why the compressed norm is exact}

The source scalar norm is the product over the main square-root conjugates
of a polynomial encoding the first, strong and auxiliary Pell conditions.
For its quadratic-remainder evaluator, the two final coefficients are
\[
 v_r=a_1(a_1\mathcal T+2b_1)-16d_nj,
 \qquad u_r=b_1^2-a_1^2\mathcal K,
\]
and the final norm is
\[
 u_r^2+\mathcal T u_rv_r+\mathcal K v_r^2.
\]
The identity
\begin{equation}
 u_r^2+\mathcal T u_rv_r+\mathcal K v_r^2
 =\mathcal K v^2-u(\mathcal T v-u),\qquad
 u=-u_r,\quad v=v_r,
 \label{eq:factored-positive-norm}
\end{equation}
holds in every commutative ring.  The sign change reverses an existing
subtraction, and the factorization saves one multiplication.

Two additional ring identities give the other savings.  Sharing $H^2Q$
between $H_P=H^2QDc^2T^2$ and $H_U=H^2Qb^2$ saves one multiplication.
Computing
\[
 d^2(W^2-L^2G)^2=[d(W^2-L^2G)]^2
\]
saves another because the source's separate $d^2$ has no other consumer.
Thus the product takes $108=65+43$ operations, the summands being
multiplications and additions/subtractions.  The final comparison in
\eqref{eq:system103} avoids one more subtraction, giving
$103=63+40$.  The transport shift gives $102=63+39$.

No new rational-root, root-separation or Pell-saturation hypothesis is
introduced by these changes.  The norm is the identical four-port
polynomial template.  Consequently the source's universal reconstruction,
its product-unit sign argument and its nonzero leading homogeneous form
remain valid.  The norm has exact degree $1072$, and $N_k,N_t$ have degrees
$7,2$, respectively, so \eqref{eq:p108} has exact degree $1081$.
The active coefficient tuple has $\beta>0$, which is the coefficient
condition needed for the inherited noncancellation certificate.

\subsection{Logical consequence and reproducibility}

For each presentation in this section, the new complete coefficient
adapter, the inherited arithmetic reconstruction, and the literal K215
calibration give
\begin{align*}
 \text{positive solvability}
 &\quad\Longleftrightarrow\quad \text{blank halting of K215}\\
 &\quad\Longleftrightarrow\quad \text{existence of a finite ZFC refutation}.
\end{align*}
Negation yields the requested equivalence with $\operatorname{Con}(\mathrm{ZFC})$
inside ZFC.  The construction establishes neither side of that equivalence.

The quartic generator performs exact sparse-polynomial verification of
all residuals and of the degree certificate.  The twelve-variable checker
verifies five local ring identities exactly, audits every live input and
operation, compares $240$ signed modular assignments with the source
template, and checks twelve additional assignments using residues of the
actual bound coefficients.  Both active bindings authenticate the same
four payloads.  Their historical source files have been restored to exact
repository bytes and checked against the pinned Git blob identifiers.
The new algebraic and coefficient arguments remain ordinary mathematical
proofs; these finite checks are not a new end-to-end Lean or Rocq theorem.

\section{Reproduction, attribution and remaining work}
\label{sec:reproduction}

\subsection{What is in the archive}

The LaTeX source consists of \texttt{report.tex}, the files under
\texttt{sections/}, and \path{signed10/SIGNED10_SECTION.tex}.
The compiled \texttt{report.pdf} contains the complete
mathematical specification.  The four implementation directories serve
distinct purposes:
\begin{description}[leftmargin=28mm,style=nextline]
 \item[\texttt{signed10/}] The ten-integer equation, four comparison
 variants, explicit fixed-leaf recipes, norm polynomials, coefficient
 payloads, a builder, and author and independent audits.
 \item[\texttt{signed16/}] The independent Jones-based sixteen-integer
 equation, its index data, coefficient recipes, circuit and verifier.
 \item[\texttt{positive12/}] The 108-operation product, 103-operation
 unshifted system and 102-operation shifted system, with symbolic and
 concrete-instance checks.  It also contains an independent audit of the
 new native coefficient construction.
 \item[\texttt{quartic31/}] The 60-operation quadratic system,
 125-operation quartic and exact symbolic coefficient data, together with
 an independent literal audit of the sixteen-variable equation.
 \item[\texttt{coefficients/}] The literal machine transitions and tile
 table, the complete finite tuple recipe, raw and compressed port bytes,
 modular/controller certificates, and the adapter proof.
\end{description}
Selected pinned source proofs are under \texttt{sources/} and the local
\texttt{upstream/} directories.  They are evidence and context; they are
not a standalone checkout of ProveIt or all of its Lean dependencies.
The original machine notices and repository license are preserved.

\subsection{Commands and verification scope}

From the extracted archive root, run
\begin{verbatim}
python3 verify_all.py
make pdf
\end{verbatim}
The arithmetic checks require Python 3.11 or later and only its standard
library.  They need no network access or solver.  Do not use Python's
\texttt{-O} option: many exact assertions are part of the audit, and the
top-level driver rejects optimized execution.  The PDF build requires a
standard pdfLaTeX installation with the packages listed in the preamble.
The scripts resolve their data relative to their own locations, so the
tree can be moved into a new numbered ProveIt report directory unchanged.

The finite audits establish the following concrete facts.  Every emitted
operation has two earlier inputs and is an addition, subtraction or
multiplication.  Every advertised variable is live.  All actual port
bytes and their hashes agree with the complete tuple.  The parent
normalization reconstructs the 402 weighted coefficient records exactly.
The index and Sun-source bounds hold for those actual integers.  The
three-root and four-root norm identities are checked as exact polynomial
identities.  Separate literal evaluators agree with the emitted circuits
on signed and boundary assignments modulo several primes.  These
assignment checks supplement the mathematical transcription proofs; they
do not replace them.

Degree verification is stronger than a collection of sample values.
Gate-by-gate propagation gives an upper bound.  Substituting a common
indeterminate into the variable leaves, while retaining exact leading
coefficients modulo a specified prime, produces a nonzero coefficient at
the claimed top degree.  Thus the bound is attained for the actual fixed
tuple.  The low-degree systems additionally have exact symbolic residual
comparisons; the quartic has the uniform coefficient $[\tau^4]Q=1$.

\subsection{What each improvement contributes}

The native tuple combines the unchanged mixed-radix tile geometry with a
different odd-factor rank controller.  Its contribution is smaller fixed
ports, not a smaller machine or a different simulated language.  The
ten-integer equation specializes an existing universal reduction to this
literal source, then reduces the number of square tests, the homogeneous
merger degree and the root-separation weights.  The sixteen-integer
equation develops a separate base restriction and guard elimination; its
large circuit improvement also uses the already established digit
evaluation identities from package 27.  The 12- and 31-variable changes
are exact algebraic refactorings and a justified positive-coordinate
shift.  Their savings should be interpreted at their stated degree and
domain, not as a new universal minimum operation count.

\subsection{Useful next steps}

The most valuable next verification step is to formalize the new
coefficient adapter, the signed norm lemma, the quintic merger and the
cofinal-base lemma, then connect their literal circuit definitions to
the existing K215 endpoint.  For the ten-variable route the modified
Lucas-certificate and residue-interpolation theorems of the 2026 preprint
also belong in that formalization boundary.

Further arithmetic sharing may reduce the 1,442-operation evaluator; no
addition-chain or circuit minimality certificate is supplied.  A smaller
native machine, a different tile language, or a smaller suitable
low-degree source could improve several metrics at once, but would
require a new calibrated compiler argument.  Reducing ten unrestricted
integer variables is a separate research problem.  The present work
does not claim that the general ten-unknowns bound, the nine-variable
natural-number bound, or this particular consistency instance is optimal.


\begin{thebibliography}{99}
\bibitem{repo}
V. Reshetnikov and contributors, \emph{ProveIt: Con(ZFC) Diophantine forms},
source snapshot \texttt{8c6eec9d19ef80538c8e61ff6e0180040583446a}.
\url{https://github.com/VladimirReshetnikov/ProveIt/tree/8c6eec9d19ef80538c8e61ff6e0180040583446a/SetTheory/Cardinals/docs/reports/hilbert-tenth-problem/con-zfc-diophantine-forms}.

\bibitem{repo24}
ProveIt, package 24, \emph{integer24}; in particular the general
odd-factor controller in its coefficient construction, at the snapshot
in \cite{repo}.  Relevant source is included under
\texttt{coefficients/upstream/24-coefficients.tex}.

\bibitem{repo25}
ProveIt, package 25, \emph{hybrid12}, mixed-radix Wang943 coefficient
construction, at the snapshot in \cite{repo}.  Included source:
\texttt{coefficients/upstream/25-COEFFICIENT\_PROOF.md}.

\bibitem{repo26}
ProveIt, package 26, \emph{ratio12}, including the signed product and
ratio normalizations, at the snapshot in \cite{repo}.

\bibitem{repo27}
ProveIt, package 27, \emph{wave4}, digit-evaluation circuits and
coefficient/cost comparisons, at the snapshot in \cite{repo}.

\bibitem{repo29}
ProveIt, package 29, \emph{simpler-forms}, including the 17-integer
Jones wrapper, at the snapshot in \cite{repo}.  Included source:
\texttt{sources/SIGNED17\_SOURCE\_PROOF.md}.

\bibitem{repo30}
ProveIt, package 30, \emph{quartic31}: the positive quartic parent,
the normalized 31-variable source, and the 12-variable specialization,
at the snapshot in \cite{repo}.

\bibitem{endpoint}
ProveIt, \texttt{KEndpoint215.lean}, calibrated blank-input K215 endpoint,
at the snapshot in \cite{repo}.  The exact source used here is bundled
as \texttt{sources/KEndpoint215.lean}.

\bibitem{jones}
J. P. Jones, Universal Diophantine equation,
\emph{Journal of Symbolic Logic} \textbf{47} (1982), 549--571.
\url{https://doi.org/10.2307/2273588}.

\bibitem{sun2021}
Z.-W. Sun, Further results on Hilbert's Tenth Problem,
\emph{Science China Mathematics} \textbf{64} (2021), 281--306.
Version used: arXiv:1704.03504v7, especially the explicit construction
in Section 3 and the square-combining polynomial in Section 5.
\url{https://arxiv.org/html/1704.03504v7}.

\bibitem{sun2026}
Z.-W. Sun, Ten unknowns for Hilbert's tenth problem over the integers,
preprint, arXiv:2609.01594v2 (2026), especially Sections 3--7.
\url{https://arxiv.org/html/2609.01594v2}.
\end{thebibliography}
\end{document}
